\section{Proof of guarantees using Tillé's pivot (Theorem~\ref{th:tille})}
\label{app:tille}

\begin{lemma} \label{lemma:top_tille}
    The top recommendation of Tillé's algorithm matches that of \SimpleVeche{},
    i.e. $\TilleVeche_{1} (b) = \SimpleVeche_{1} (b)$ 
    for all ballots $b \in \setR^{U \times E}$.
\end{lemma}

\begin{proof}
    We prove this by induction over the number $E$ of entities.
    If $E = 2$, then it is trivial.

    Assume it holds for $E$ entities.
    Consider $E + 1$ entities,
    and let $b \in \setR^{U \times E}$ and denote $w = \Weights(b)$.
    Let $e^* = \TilleVeche_1 (b)$ be the random top entity,
    Let $(e_1, e_2) \in [E + 1]^2$ be the random first pair of entities 
    to be opposed in the first round, with $e_1 \neq e_2$,
    and $x \in \set{1, 2}$ be the result of their opposition 
    (with $x = 1$ if $e_1$ wins).
    By the law of total probability, we have
    \begin{align}
        \pb{e^* = e|b}
        &= \bbE_{e_1,e_2} \left[
            \bbE_x \left[ \pb{e^* = e|w, e_1, e_2, x} \right]
        \right] \\
        &= \bbE_{e_1,e_2} \left[
            \frac{
                w_{e_1} \pb{e^* = e | w^{e_1 > e_2}}
                + w_{e_2} \pb{e^* = e | w^{e_2 > e_1}}
            }{w_{e_1} + w_{e_2}}
        \right], \label{eq:total_law}
    \end{align}
    where $w^{e_1 > e_2}$ is the resulting weight vector
    after $e_1$ stealing the weights of $e_2$.
    By the induction assumption,
    the distribution of $e^*$ given $w^{e_1 > e_2}$
    matches sampling without replacement,
    which leads to
    \begin{align}
        \pb{e^* = e | w^{e_1 > e_2}}
        = \frac{w^{e_1 > e_2}_{e}}{\norm{w^{e_1 > e_2}}{1}}
        = \frac{w^{e_1 > e_2}_{e}}{\norm{w}{1}},
    \end{align}
    using the fact that the weights have been transfered but not ``lost''.
    For all $e \notin \set{e_1, e_2}$,
    we then have $\pb{e^* = e | w^{e_1 > e_2}} = w_e / \norm{w}{1}$.
    Moreover, $\pb{e^* = e_1 | w^{e_1 > e_2}} = (w_{e_1} + w_{e_2}) / \norm{w}{1}$
    and $\pb{e^* = e_2 | w^{e_1 > e_2}} = 0$.
    Overall, for all $e \in [E+1]$, we thus have
    \begin{align}
        \frac{
            w_{e_1} \pb{e^* = e | w^{e_1 > e_2}}
            + w_{e_2} \pb{e^* = e | w^{e_2 > e_1}}
        }{w_{e_1} + w_{e_2}}
        = \frac{w_e}{\norm{w}{1}}.
    \end{align}
    Plugging this into the total law equation~\eqref{eq:total_law}
    yields the lemma.
\end{proof}

\begin{proof}[Proof of Theorem~\ref{th:tille}]
    This follows straightforwardly 
    from Theorem~\ref{th:simpleveche} and Lemma~\ref{lemma:top_tille}.
\end{proof}


% \subsection{Similarity guarantees}

% \begin{proof}
%     If $e$ and $f$ were posted by the same councillor $u$,
%     then $b_{ue}, b_{uf} > 0$.
%     If no other councillor interacted with any, then $b_{ve}, b_{vf} = 0$.
%     Then $b_{\cdot e}$ and $b_{\cdot f}$ are clearly colinear 
%     and in the same direction (namely that of the $u$-th canonical basis vector),
%     hence $\rho_{ef} = 1$.

%     If $e$ and $f$ were posted by two different councillors $u$ and $v$,
%     and if no other councillor interacted with any, 
%     then $b_{\cdot e}$ is colinear with the $u$-th canonical basis vector,
%     while $b_{\cdot f}$ is colinear with the $v$-th canonical basis vector.
%     They are orthogonal, hence $\rho_{ef} = 0$.

%     In the last case, $b_{\cdot e}$ and $b_{\cdot f}$ point in opposite direction,
%     along the vector ${\bf e}_u - {\bf e}_v$ 
%     (where these correspond to canonical basis vector).
% \end{proof}
